\documentclass[14pt,a4paper]{extarticle}
\usepackage[T2A]{fontenc}
\usepackage[utf8]{inputenc}
\usepackage[english,russian]{babel}
\usepackage{lmodern}
\usepackage{geometry}
\usepackage{array}
\usepackage{tabularx}
\usepackage{longtable}
\usepackage{ltablex}
\usepackage{booktabs}
\usepackage{hyperref}
\usepackage{indentfirst}
\usepackage{microtype}
\usepackage{graphicx}
\usepackage{setspace}
\usepackage{listings}
\usepackage{graphicx}
\usepackage{caption}
\usepackage{float}
\usepackage{tocloft}
\usepackage{titlesec}
\usepackage{newtxtext,newtxmath}

\titleformat{\section}        % форматируем \section
  {\normalfont\bfseries}      % стиль: обычный + жирный
  {\thesection.}              % номер секции с точкой
  {2mm}                       % расстояние между номером и текстом
  {}                           % перед заголовком можно добавить что-то
\titlespacing*{\section}{1.5cm}{*3}{*1}  % {отступ слева}{отступ сверху}{отступ снизу}

\titleformat{\subsection}        % форматируем \section
  {\normalfont\bfseries}      % стиль: обычный + жирный
  {\thesubsection.}              % номер секции с точкой
  {2mm}                       % расстояние между номером и текстом
  {}                           % перед заголовком можно добавить что-то
\titlespacing*{\subsection}{1.5cm}{*3}{*1}  % {отступ слева}{отступ сверху}{отступ снизу}

\captionsetup{justification=centering}
\renewcommand{\ttdefault}{cmtt}
\keepXColumns
\setlength{\parindent}{3em}
\newgeometry{left=2.5cm,right=2.5cm,top=3cm,bottom=2.5cm}

\renewcommand{\contentsname}{\hspace*{4em}{\smallСодержание}}
\setlength{\cftbeforesecskip}{0.5em}
\renewcommand{\cftsecleader}{\cftdotfill{\cftdotsep}}
\renewcommand{\cftsecaftersnum}{.\hspace{2mm}}

\title{Отчёт по "Управление ИТ-проектами"}
\author{Выполнил студент группы М-РИТ-25-1, Станиславчук Сергей}

\begin{document}

\setstretch{1.5}
\noindent
\begin{minipage}[t]{0.01\textwidth}
\includegraphics{/home/stanik/programmer/++Programmer/5_2/LSTU.jpg}
\end{minipage}%
\begin{minipage}[c]{1.1\linewidth}
    \centering
    {\small
    МИНИСТЕРСТВО НАУКИ И ВЫСШЕГО ОБРАЗОВАНИЯ РФ \\[1em]
    ФЕДЕРАЛЬНОЕ ГОСУДАРСТВЕННОЕ БЮДЖЕТНОЕ \\
    ОБРАЗОВАТЕЛЬНОЕ УЧРЕЖДЕНИЕ ВЫСШЕГО ОБРАЗОВАНИЯ \\
    «ЛИПЕЦКИЙ ГОСУДАРСТВЕННЫЙ ТЕХНИЧЕСКИЙ УНИВЕРСИТЕТ»
    }
\end{minipage}%

\vspace{3em}

\noindent
\hspace{1em}%
\begin{minipage}[t]{0.8\textwidth}
\textbf{Институт} \hspace{2em} компьютерных наук \\
\textbf{Кафедра} \hspace{2em} автоматизированных систем управления
\end{minipage}


\vspace{6em}

\begin{center}
\textbf{Практическая работа №14} \\
по дисциплине \\ <<Управление ИТ-проектами>>
\end{center}

\vspace{6em}

\noindent
\begin{tabular*}{\textwidth}{@{\extracolsep{\fill}} l c r }
    Студент М-РИТ-25-1 & \makebox[4cm]{\hrulefill} & Станиславчук С.М. \\[2em]
Руководитель & \\
К.т.н., доцент & \makebox[4cm]{\hrulefill} & Ведищев В.В. \\
\end{tabular*}

\vfill
\begin{center}
    Липецк, 2026
\end{center}

\thispagestyle{empty}

\newpage
\tableofcontents
\newpage

\section{Составление плана управления качеством программного продукта}

\subsection{Ответственность за качество. Контроль качества}

Обеспечение качества программного продукта является обязанностью всей команды разработки, однако степень ответственности участников различается в зависимости от выполняемых функций. В рамках разрабатываемой системы корпоративного обмена сообщениями бэкенд-разработчик отвечает за корректность серверной логики, обработку сетевых запросов, устойчивость взаимодействия компонентов и соблюдение требований к производительности. Фронтенд-разработчик обеспечивает качество пользовательского интерфейса, удобство взаимодействия с системой и корректность отображения данных.

Специалист по информационной безопасности осуществляет контроль реализации механизмов защиты, включая аутентификацию, шифрование данных и управление доступом. Тестировщик отвечает за выявление дефектов, проведение функционального и регрессионного тестирования, а системный администратор контролирует корректность развёртывания, стабильность инфраструктуры и эксплуатационную готовность системы.

Контроль качества осуществляется на протяжении всего жизненного цикла разработки и включает анализ требований, проверку архитектурных решений, инспектирование исходного кода, тестирование и контроль стабильности релизов. Для снижения вероятности ошибок используется система контроля версий Git, а также процедуры проверки изменений перед включением в основную ветку проекта.

\subsection{Подготовка и проведение инспектирования}

Инспектирование представляет собой процедуру проверки промежуточных результатов разработки с целью раннего выявления ошибок и отклонений от требований. В проекте инспектирование проводится в отношении исходного кода, архитектурных решений, документации и конфигурационных файлов.

Подготовка инспектирования включает определение объекта проверки, перечня критериев качества и ответственных лиц. Например, при инспектировании серверной части оценивается корректность сетевого взаимодействия, обработка ошибок, безопасность обмена сообщениями и соответствие требованиям производительности.

Проведение инспектирования осуществляется в форме анализа изменений (code review), обсуждения архитектурных решений и проверки документации. По результатам инспектирования фиксируются замечания, после чего выполняется корректировка выявленных недостатков.

\subsection{План контроля качества – состав}

План контроля качества определяет перечень мероприятий, обеспечивающих соответствие системы установленным требованиям. В состав плана контроля качества входят:

\begin{itemize}
    \item контроль соответствия функциональных возможностей требованиям спецификации;
    \item проверка корректности реализации механизмов аутентификации и авторизации;
    \item тестирование передачи текстовых сообщений, файлов и голосовых данных;
    \item нагрузочное тестирование системы при одновременной работе нескольких пользователей;
    \item проверка устойчивости к сетевым сбоям и восстановлению соединения;
    \item анализ безопасности передачи данных и корректности механизмов шифрования;
    \item контроль стабильности сборки и корректности развёртывания релизов;
    \item анализ журналов работы системы и мониторинг ошибок.
\end{itemize}

Контроль качества выполняется итеративно на каждом этапе разработки, что позволяет снизить стоимость исправления ошибок и повысить устойчивость системы.

\subsection{Верификация}

Верификация представляет собой процесс подтверждения того, что программный продукт реализован в соответствии с установленными требованиями и проектной документацией. В контексте разрабатываемой системы верификация направлена на проверку корректности реализации функциональных требований, включая авторизацию пользователей, обмен сообщениями, голосовую связь и передачу файлов.

Верификация выполняется посредством анализа исходного кода, модульного тестирования, интеграционного тестирования и проверки соответствия результатов разработки технической документации. Например, производится проверка того, что сообщения доставляются адресату, а механизмы аутентификации работают в соответствии с заданными правилами.

Таким образом, верификация отвечает на вопрос: \textit{«Правильно ли реализована система?»}

\subsection{Валидация}

Валидация представляет собой процесс подтверждения того, что разработанная система удовлетворяет реальным потребностям пользователей и пригодна для эксплуатации в корпоративной среде. В отличие от верификации, валидация ориентирована на оценку практической применимости системы.

В рамках проекта валидация включает тестирование системы в условиях, приближенных к реальной корпоративной сети, оценку удобства интерфейса, качества голосовой связи, скорости доставки сообщений и устойчивости к нагрузкам.

Проведение валидации позволяет определить, удовлетворяет ли система ожиданиям пользователей и способна ли эффективно использоваться для корпоративного обмена сообщениями.

Таким образом, валидация отвечает на вопрос: \textit{«Правильную ли систему разработали?»}

\end{document}
